\newpage

\renewcommand\thesection{2}
\renewcommand\thesubsection{\arabic{subsection}}

\counterwithin{subsection}{section}
\setcounter{subsection}{0}

\begin{center}
    \section*{\textbf{BAB II \\ SOLUSI PERMASALAHAN}}
    \addcontentsline{toc}{section}{BAB II \quad SOLUSI PERMASALAHAN}
\end{center}

\subsection{Solusi 1}
    Solusi pertama yang dapat dilakukan untuk menjawab permasalahan mengenai kesenjangan antara pengetahuan agama dan perilaku antikorupsi adalah dengan memperkuat pembinaan moral yang lebih kontekstual. Selama ini, nilai agama sering kali dipahami secara normatif, seperti mengetahui bahwa korupsi adalah tindakan yang dilarang, tetapi belum sepenuhnya diterapkan dalam kehidupan nyata. Oleh karena itu, pembinaan keagamaan perlu diarahkan tidak hanya pada penyampaian doktrin, tetapi juga pada pembahasan kasus nyata yang berhubungan dengan dilema moral dalam kehidupan sehari-hari.

    Diskusi moral yang kontekstual dapat dilakukan melalui kajian, ceramah, seminar, atau kegiatan pembinaan yang membahas contoh konkret seperti suap, gratifikasi, penyalahgunaan jabatan, manipulasi data, mencontek, titip absen, dan perilaku tidak jujur lainnya. Melalui pendekatan ini, masyarakat tidak hanya mengetahui bahwa korupsi merupakan tindakan yang salah, tetapi juga memahami bagaimana cara menolak dan mencegah perilaku koruptif ketika menghadapi tekanan sosial atau kepentingan pribadi. Pembinaan seperti ini diharapkan dapat membantu masyarakat menghubungkan nilai agama dengan tindakan nyata dalam kehidupan sehari-hari \citep{hasan2024,nuruddin2024}.


\subsection{Solusi 2}  
     Solusi kedua adalah meningkatkan peran tokoh agama sebagai figur teladan dalam membangun budaya antikorupsi. Tokoh agama memiliki pengaruh yang besar karena dipandang sebagai panutan moral oleh umat dan masyarakat. Oleh sebab itu, tokoh agama tidak hanya berperan dalam menyampaikan ajaran agama dari mimbar, tetapi juga perlu hadir sebagai contoh nyata dalam menerapkan nilai kejujuran, tanggung jawab, kesederhanaan, dan integritas di ruang publik.

     Keteladanan tokoh agama menjadi penting karena masyarakat membutuhkan figur yang dapat menunjukkan bahwa nilai agama tidak berhenti pada kegiatan ritual, tetapi juga diwujudkan dalam sikap dan keputusan sehari-hari. Tokoh agama dapat berperan aktif dengan menyuarakan pesan antikorupsi dalam ceramah, membahas isu integritas dalam kegiatan keagamaan, serta mendorong umat untuk berani menolak segala bentuk tindakan tidak jujur. Dengan adanya figur teladan yang konsisten, nilai antikorupsi dapat lebih mudah diterima dan diterapkan oleh masyarakat, khususnya generasi muda \citep{siagian2025,zahron2026}.

        
\subsection{Solusi 3}
    Solusi ketiga adalah mengintegrasikan pendidikan antikorupsi berbasis nilai agama dalam lingkungan pendidikan dan masyarakat. Pendidikan antikorupsi perlu diberikan sejak dini agar nilai kejujuran, tanggung jawab, disiplin, dan kepedulian terhadap hak orang lain dapat terbentuk sebagai kebiasaan. Dalam lingkungan pendidikan, nilai antikorupsi dapat dikaitkan dengan perilaku sederhana seperti tidak mencontek, tidak melakukan plagiarisme, tidak memanipulasi data, dan berani bertanggung jawab atas tindakan sendiri.

    Integrasi pendidikan antikorupsi juga perlu melibatkan keluarga, sekolah, lembaga keagamaan, dan masyarakat. Keluarga berperan dalam membiasakan kejujuran sejak kecil, sekolah berperan dalam membentuk karakter melalui aturan dan pembelajaran, sedangkan lembaga keagamaan berperan dalam memberikan dasar moral dan spiritual. Melalui kerja sama berbagai pihak, nilai agama dapat diterapkan secara lebih nyata dalam kehidupan sosial, sehingga masyarakat tidak hanya memahami bahwa korupsi adalah tindakan yang salah, tetapi juga memiliki keberanian moral untuk menolaknya \citep{devina2026,nuruddin2024}.
